\documentclass{article}
\usepackage{indentfirst}
\usepackage{listings}
\usepackage{amsmath}
\usepackage{amssymb}
\usepackage{pifont}
\usepackage{url}
\usepackage{tikz-cd}
\usepackage[utf8]{inputenc}
\usepackage[english,ukrainian]{babel}
\usepackage{booktabs}
\usepackage{hyperref}
\usepackage{geometry}
\geometry{a4paper, margin=2.5cm}
\input{journal}

\hypersetup{
    colorlinks=true,
    linkcolor=blue,
    citecolor=blue,
    urlcolor=blue,
    filecolor=magenta
}

\lstset{
  basicstyle=\ttfamily\small,
  breaklines=true,
  frame=single,
  numbers=left,
  numberstyle=\tiny\color{gray},
  keywordstyle=\color{blue},
  commentstyle=\color{green!60!black},
  stringstyle=\color{red},
  showstringspaces=false,
  language=,
  tabsize=2
}

\begin{document}

\title{ПОРІВНЯЛЬНИЙ АНАЛІЗ ІНСТРУМЕНТАРІЮ ГЕНЕРАЦІЇ ДОКУМЕНТАЦІЇ У ФОРМАТІ API BLUEPRINT ПРИ МІГРАЦІЇ З ORACLE APIARY: \\ КЕЙС СИСТЕМИ ЕСОЗ УКРАЇНИ}
\author{Максим Сохацький$^1$}
\date{%
$^1$ Керівник відділу системного аналізу ДП «Електронне Здоров'я»\\[0.3em]
Національний технічний університет України \\
\small Ігор Сікорський Київський політехнічний інститут \\
\today
}
\maketitle

УДК 004.415.2:614.2(477)

SCOPUS: \url{https://www.scopus.com/authid/detail.uri?authorId=57214352719}

ORCID: \url{https://orcid.org/0000-0001-7127-8796}

\begin{abstract}
У статті проведено порівняльний аналіз чотирьох активно підтримуваних інструментів рендерингу API-документації у форматі Apiary 1A (API Blueprint): \textbf{Blueprinter}, \textbf{Aglio}, \textbf{Api-blueprint-boilerplate} та \textbf{Phpdraft}. Дослідження зумовлене необхідністю міграції документації Єдиної системи охорони здоров'я України (ЕСОЗ) у зв'язку із завершенням строку підтримки (End-of-Life, EoL) платформи Oracle Apiary. На підставі формалізованого набору критеріїв — повноти підтримки MSON, якості обробки попереджень парсера, зручності УІ, автономності від хмарних сервісів Oracle та придатності до роботи з масштабними специфікаціями — обґрунтовано вибір \textbf{Blueprinter} як оптимального інструменту для безпомилкової міграції. Практичним артефактом дослідження є репозиторій \texttt{ehealth-ua/ONE-API} та мігрований файл \texttt{esoz-blueprint.apib}.

\textbf{Ключові слова:} API Blueprint, APIB, Oracle Apiary, ЕСОЗ, міграція документації, Blueprinter, Aglio, MSON, рендерер, сумісність.
\end{abstract}

\begin{abstract}
This paper presents a comparative analysis of four actively maintained API Blueprint (Apiary 1A format) documentation rendering tools: Blueprinter, Aglio, Api-blueprint-boilerplate, and Phpdraft. The study is motivated by the necessity of migrating the Ukrainian Electronic Health System (eHealth) API documentation following the End-of-Life announcement for Oracle Apiary. Based on a formalized set of criteria — MSON support completeness, parser warning handling quality, UI ergonomics, independence from Oracle cloud services, and scalability for large specifications — \textbf{Blueprinter} is justified as the optimal migration target. The practical artifact of this research is the \texttt{ehealth-ua/ONE-API} repository and the migrated \texttt{esoz-blueprint.apib} specification file.

\textbf{Keywords:} API Blueprint, APIB, Oracle Apiary, eHealth Ukraine, documentation migration, Blueprinter, Aglio, MSON, renderer, compatibility.
\end{abstract}

\section{Вступ}
Документація програмних інтерфейсів (API) є стратегічним компонентом архітектури сучасних державних інформаційних систем. У рамках цифрової трансформації охорони здоров'я України функціонує Єдина система охорони здоров'я (ЕСОЗ) — централізована інформаційна платформа, що інтегрує медичні інформаційні системи (МІС), лабораторії, аптеки та регуляторні органи. Взаємодія між учасниками екосистеми здійснюється через REST API, специфікації якого традиційно описуються у форматі \textbf{API Blueprint} і публікуються через платформу \textbf{Oracle Apiary}.

Проте Oracle оголосила про завершення активної підтримки платформи Apiary (End-of-Life), що ставить перед командою ЕСОЗ завдання вибору альтернативного інструментарію. Критичним є збереження можливості рендерингу існуючих специфікацій формату APIB без значних переробок, оскільки основний файл \texttt{apiary.apib} має об'єм понад 1 МБ і охоплює сотні ресурсів, типів даних MSON та унікальних розширень формату.

\textbf{Мета статті} — провести систематичний порівняльний аналіз чотирьох альтернативних інструментів та науково обґрунтувати доцільність вибору Blueprinter як основного рендерера для документації ЕСОЗ.

\section{Огляд формату API Blueprint та його інструментальної екосистеми}
API Blueprint — це мова опису веб-API рівня документації, розроблена компанією Apiary (Jakub Nešetřil та Zdenko Vrábel, 2013). Вона ґрунтується на синтаксисі Markdown та MSON (Markdown Syntax for Object Notation) і зберігається у файлах з розширенням \texttt{.apib} \cite{apiblueprint_spec}. Специфікація підтримується відкритою бібліотекою \texttt{drafter} (C++) та парсером \texttt{minim} (JavaScript), які реалізують перетворення тексту у дерево API Elements (IETF Internet Draft).

ЕСОЗ з самого початку побудована навколо цього формату. Всі специфікації REST-ресурсів — від \texttt{Patient}, \texttt{Encounter}, \texttt{Composition} до \texttt{MedicationRequest} — описані у \texttt{.apib} файлах та зводяться в єдину специфікацію \texttt{apiary.apib} обсягом понад мільйон байт.

Ключовою відмінністю між доступними рендерерами є глибина їхньої інтеграції з виходом парсера: деякі з них працюють лише з AST-представленням першого покоління (\texttt{apib-ast}), тоді як сучасніші — з API Elements другого покоління, що є більш семантично насиченим.

\section{Формальний аналіз порівнюваних інструментів}

\subsection{Aglio}
\textbf{Автор та генеалогія.} Daniel G. Taylor (GitHub: \texttt{danielgtaylor}). Проєкт \texttt{danielgtaylor/aglio} — один із перших і найбільш відомих рендерерів API Blueprint. Початково презентований близько 2014 року, він дав поштовх розвитку всієї екосистеми відкритих рендерерів APIB. Має понад 4 000 зірок на GitHub.

\textbf{Технічні характеристики.} Написаний на Node.js. Ядро (core) завантажує \texttt{.apib}-файл та обробляє аргументи, а тематичний рушій (theme engine) конвертує AST у HTML. За замовчуванням використовує тему Olio; підтримує також Streak, Flatly, Slate, Cyborg. Підтримує вбудовані директиви \texttt{<!-- include -->} для поєднання кількох APIB-файлів.

\textbf{Сильні сторони:}
\begin{itemize}
\item Велика база встановлень та документації.
\item Підтримка кастомних Jade-шаблонів.
\item Вбудований live-preview сервер.
\item Широка інтеграція у CI/CD пайплайни (Docker-образи, плагіни для Foliant, mkdocs тощо).
\end{itemize}

\textbf{Слабкі сторони:}
\begin{itemize}
\item З 2019 року автор офіційно повідомив про обмежену підтримку та запросив нових мейнтейнерів \cite{aglio_maintenance}.
\item Залежить від застарілих версій \texttt{drafter}/\texttt{protagonist}, що обумовлює некоректну обробку новітніх розширень MSON (зокрема, вкладених членів підтипів примітивних типів), генеруючи численні попередження при рендерингу специфікацій ЕСОЗ.
\item Відсутність нативного UX для мобільних пристроїв та темної теми.
\end{itemize}

\subsection{Api-blueprint-boilerplate}
\textbf{Автор та генеалогія.} Jakub Synowiec (GitHub: \texttt{jsynowiec}), проєкт \texttt{jsynowiec/api-blueprint-boilerplate}. Не є самостійним рендерером; являє собою готовий starter-kit (шаблон репозиторію) для швидкого старту проєктів документації у форматі APIB.

\textbf{Технічні характеристики.} Включає визначений набір залежностей: \texttt{apiary-cli} для локального перегляду та валідації (використовує хмарний API Apiary), \texttt{drafter.js} для тестування схеми, \texttt{aglio} для генерації статичного HTML. Вимагає наявності Node.js, NPM, Ruby та Bundler.

\textbf{Сильні сторони:}
\begin{itemize}
\item Готова структура проєкту з прикладами, CI-конфігурацією та документацією.
\item Вбудований \texttt{npm test} для швидкої валідації схеми через drafter.js.
\end{itemize}

\textbf{Слабкі сторони:}
\begin{itemize}
\item Принципово залежить від \texttt{apiary-cli} — компонента екосистеми Oracle Apiary, що підлягає міграції. Після EoL Apiary інструмент у базовій конфігурації стає непридатним.
\item На відміну від інших кандидатів, не є рендерером, а лише обгорткою, тому не може безпосередньо замінити Apiary у виробничому середовищі.
\item Відсутня власна система тем та функції пошуку.
\end{itemize}

\subsection{Phpdraft}
\textbf{Автор та генеалогія.} Sean Miller (GitHub: \texttt{SMillerDev}), репозиторій \texttt{SMillerDev/phpdraft}. Первинний автор — Ryan Hellenius (попередня версія). Це єдиний PHP-реалізований парсер та рендерер API Blueprint у наведеному переліку.

\textbf{Технічні характеристики.} Написаний на PHP 8.1+. Для парсингу \texttt{.apib} потребує системного бінарного файлу \texttt{drafter} або може використовувати онлайн-парсер Apiary за відсутності локального. Підтримує генерацію HTML та конвертацію у формат OpenAPI (Swagger).

\textbf{Сильні сторони:}
\begin{itemize}
\item Можливість нативної інтеграції в PHP-додатки (Laravel, Symfony).
\item Функція конвертації APIB → OpenAPI — цінна при стратегії повноцінної міграції на OpenAPI у майбутньому.
\item Підтримка CSS та JavaScript-кастомізації шаблонів.
\end{itemize}

\textbf{Слабкі сторони:}
\begin{itemize}
\item Потребує налаштування окремого середовища (PHP + Drafter), що ускладнює DevOps для команд, що використовують Node.js.
\item Залежність від системного \texttt{drafter} робить інструмент чутливим до версій бінарного файлу.
\item Менш ергономічний UX при роботі з масштабними специфікаціями: відсутність вбудованого пошуку та навігаційних фільтрів.
\end{itemize}

\subsection{Blueprinter}
\textbf{Автор та генеалогія.} Команда Funbox (GitHub: \texttt{@funbox}), репозиторій \texttt{funbox/blueprinter}. Розроблений як сучасна альтернатива Aglio з акцентом на глибоку інтеграцію з API Elements другого покоління.

\textbf{Технічні характеристики.} Написаний на Node.js. Встановлюється через npm як \texttt{@funbox/blueprinter}. Приймає \texttt{.apib}-файл та генерує самодостатній HTML-файл (\texttt{index.html}) без зовнішніх залежностей. Підтримує режим watch/live-server.

\textbf{Сильні сторони:}
\begin{itemize}
\item \textbf{Глибока інтеграція з API Elements:} використовує сучасну версію \texttt{minim} та \texttt{drafter}, що забезпечує коректний парсинг новітніх конструкцій MSON.
\item \textbf{Детальна обробка помилок парсера:} чітке відображення попереджень безпосередньо в HTML-інтерфейсі з прив'язкою до рядків файлу.
\item \textbf{Повнотекстовий пошук:} можливість пошуку по групах, ресурсах та діях.
\item \textbf{Інтерактивні «One Of» блоки:} динамічна зміна варіантів відповіді без перезавантаження сторінки.
\item \textbf{Версія для друку:} адаптований режим без зайвих елементів.
\item \textbf{Локалізація UI:} можливість задати мову інтерфейсу.
\item \textbf{Повна автономність:} не потребує сервісів Oracle; працює локально або в CI/CD.
\end{itemize}

\textbf{Слабкі сторони:}
\begin{itemize}
\item Менша спільнота порівняно з Aglio.
\item Відсутність підтримки директиви \texttt{<!-- include -->} для розбиття специфікації на файли (однак це особливість формату Bootstrap-файлів, а не критична вада).
\end{itemize}

\section{Формалізація вимог до міграції}
Виходячи з профілю ЕСОЗ (масштабна APIB-специфікація, екосистема Node.js, необхідність автономності, велика стороння аудиторія МІС-розробників), сформовано такі критерії відбору:

\begin{table}[htbp]
\centering
\caption{Критерії відбору інструменту}
\footnotesize
\begin{tabular}{@{} l p{5.0cm} p{8.5cm} @{}}
\toprule
\textbf{№} & \textbf{Критерій} & \textbf{Обґрунтування} \\
\midrule
K1 & Повна підтримка MSON 2.x & Специфікація ЕСОЗ використовує складне успадкування та вкладені елементи \\
K2 & Мінімальна кількість попереджень & Безпосередній метрик якості міграції \\
K3 & Автономність від Oracle Apiary & Ключова передумова задачі (EoL) \\
K4 & Пошук та навігація в UX & МІС-розробники потребують зручного доступу до специфікацій \\
K5 & Підтримка середовища Node.js & Відповідність поточному технологічному стеку команди \\
K6 & Активна підтримка та оновлення & Зниження ризику повторного EoL у середньостроковій перспективі \\
\bottomrule
\end{tabular}
\end{table}

\section{Результати порівняльного аналізу}

\subsection{Зведена таблиця характеристик}
\begin{table}[htbp]
\centering
\caption{Порівняння інструментів рендерингу API Blueprint}
\footnotesize
\setlength{\tabcolsep}{5pt}
\begin{tabular}{@{} p{3.8cm} p{2.4cm} p{2.4cm} p{2.0cm} p{2.4cm} @{}}
\toprule
\textbf{Характеристика} & \textbf{Aglio} & \textbf{Boilerplate} & \textbf{Phpdraft} & \textbf{Blueprinter} \\
\midrule
Мова реалізації & Node.js & Node.js/Ruby & PHP & \textbf{Node.js} \\
Початок розробки & 2014 & 2015 & 2016 & $\sim$2019 \\
Статус підтримки & Maintenance & Залежить від Apiary & Активний & \textbf{Активний} \\
Версія drafter & protagonist & drafter.js & Системний & \textbf{API Elements} \\
MSON вкладеність & Часткова & Часткова & Часткова & \textbf{Повна} \\
Попередження ЕСОЗ & \textbf{Критичні} & Критичні & Середні & \textbf{Мінімальні} \\
Повнотекст.\ пошук & --- & --- & --- & \textbf{\ding{51}} \\
Автономність & \ding{51} & \textbf{\ding{55}} & \ding{51} & \textbf{\ding{51}} \\
Версія для друку & --- & --- & Базова & \textbf{\ding{51}} \\
Локалізація UI & --- & --- & --- & \textbf{\ding{51}} \\
Npm-пакет & \texttt{aglio} & --- & Composer & \texttt{blueprinter} \\
\bottomrule
\end{tabular}
\end{table}

\subsection{Оцінювання за критеріями міграції}
\begin{table}[htbp]
\centering
\caption{Оцінювання за критеріями міграції (шкала 1--5)}
\footnotesize
\setlength{\tabcolsep}{6pt}
\begin{tabular}{@{} p{5.5cm} c c c c @{}}
\toprule
\textbf{Критерій} & \textbf{Aglio} & \textbf{Boilerplate} & \textbf{Phpdraft} & \textbf{Blueprinter} \\
\midrule
K1: Підтримка MSON 2.x & 2/5 & 2/5 & 3/5 & \textbf{5/5} \\
K2: Мінімум помилок & 2/5 & 1/5 & 3/5 & \textbf{4/5} \\
K3: Автономність від Oracle & 4/5 & 1/5 & 4/5 & \textbf{5/5} \\
K4: Пошук та навігація & 2/5 & 1/5 & 2/5 & \textbf{5/5} \\
K5: Node.js-сумісність & 5/5 & 3/5 & 1/5 & \textbf{5/5} \\
K6: Активна підтримка & 2/5 & 2/5 & 4/5 & \textbf{5/5} \\
\midrule
\textbf{Загальний бал} & \textbf{17/30} & \textbf{10/30} & \textbf{17/30} & \textbf{29/30} \\
\bottomrule
\end{tabular}
\end{table}

\section{Практична апробація: міграція \texttt{apiary.apib} $\rightarrow$ \texttt{esoz-blueprint.apib}}
Для верифікації теоретичних висновків було проведено практичну міграцію у репозиторії \texttt{ehealth-ua/ONE-API}. Оригінальний файл \texttt{apiary.apib} (1 061 617 байт) було адаптовано до вимог Blueprinter та збережено як \texttt{esoz-blueprint.apib} (1 007 359 байт). Серед виконаних коригувань:
\begin{enumerate}
\item \textbf{Усунення конструкцій \texttt{sub-types of primitive types should not have nested members}} — системне попередження, яке Blueprinter відображає явно, а Aglio ігнорує, проте генерує некоректний HTML.
\item \textbf{Виправлення порожніх \texttt{Enum} секцій} — ідентифіковано та виправлено помилки \texttt{Enum element should include members}, що призводили до збоїв рендерингу.
\item \textbf{Нормалізація URI-шаблонів} — усунення дублікатів та неканонічних параметрів, що спричиняли попередження парсера.
\item \textbf{Виправлення посилань \texttt{Include}} — видалено некоректні \texttt{Include}-директиви в нетипових позиціях.
\end{enumerate}

Результуючий файл \texttt{esoz-blueprint.apib} успішно рендериться командою:
\begin{lstlisting}[language=bash]
blueprinter -i esoz-blueprint.apib -o index.html
\end{lstlisting}
без критичних помилок та з мінімальною кількістю несуттєвих попереджень (детальний звіт завжди у файлі warnings.txt). Генерований \texttt{index.html} (10 716 937 байт) є автодостатнім документом, доступним для перегляду у будь-якому сучасному браузері без підключення до мережі.

Для зручності розробки у репозиторії визначено npm-скрипти:
\begin{lstlisting}[language={}]
{
  "scripts": {
    "dev": "blueprinter -i apiary.apib -s -p 3000",
    "doc": "blueprinter -i apiary.apib -o index.html"
  },
  "dependencies": {
    "@funboxteam/blueprinter": "^6.1.0"
  }
}
\end{lstlisting}
Ці скрипти забезпечують як розробницький режим з hot-reload, так і генерацію фінального артефакту документації.

\section{Висновки}
За результатами дослідження встановлено:
\begin{enumerate}
\item \textbf{Aglio}, незважаючи на широку популяризованість, перебуває у стані обмеженої підтримки та має системні проблеми з парсингом сучасних конструкцій MSON, що критично важливо для ЕСОЗ.
\item \textbf{Api-blueprint-boilerplate} є стартовим шаблоном, а не повноцінним продуктом, і принципово залежить від хмарних сервісів Oracle Apiary, що унеможливлює його розгляд як цільової платформи після EoL.
\item \textbf{Phpdraft} є функціонально придатним варіантом, однак потребує PHP-інфраструктури та не забезпечує належного UX для широкої аудиторії МІС-розробників.
\item \textbf{Blueprinter} демонструє найвищу повноту підтримки формату API Blueprint 1A, мінімальну кількість помилок при рендерингу, сучасний та ергономічний UI, повну автономність від Oracle Apiary та відповідність технологічному стеку команди ЕСОЗ.
\end{enumerate}

Практична апробація підтвердила архітектурну правильність вибору: репозиторій \texttt{ehealth-ua/ONE-API} з мігрованим файлом \texttt{esoz-blueprint.apib} є відтворюваним науковим артефактом, що документує весь процес міграції. Blueprinter рекомендується як стандартний інструмент рендерингу API-документації ЕСОЗ на заміну Oracle Apiary.

\section*{Список використаних джерел}
\begin{enumerate}
\item API Blueprint Specification. \url{https://apiblueprint.org/documentation/specification.html}
\item danielgtaylor/aglio: Maintenance notice. GitHub Issues. \url{https://github.com/danielgtaylor/aglio/issues}
\item funbox/blueprinter. GitHub Repository. \url{https://github.com/funbox/blueprinter}
\item SMillerDev/phpdraft. GitHub Repository. \url{https://github.com/SMillerDev/phpdraft}
\item jsynowiec/api-blueprint-boilerplate. GitHub Repository. \url{https://github.com/jsynowiec/api-blueprint-boilerplate}
\item Nešetřil, J., Vrábel, Z. API Blueprint: A Web API Documentation Language. 2013. \url{https://apiblueprint.org}
\item MSON (Markdown Syntax for Object Notation) Specification. \url{https://github.com/apiaryio/mson}
\item Ukrainian ESOZ HTTP API Repository. \url{https://github.com/ehealth-ua/esoz}
\item Oracle Apiary End-of-Life Announcement. Oracle Cloud Documentation, 2024.
\item Drafter — C++ API Blueprint Parser. \url{https://github.com/apiaryio/drafter}
\end{enumerate}

\end{document}
